% \begin{figure}[!ht]
%     \centering
%     \scriptsize
%     \begin{tikzpicture}[
%         node distance=0.5cm,
%         codebox/.style={rectangle, rounded corners=3pt, draw=black!40, line width=0.8pt, 
%                         fill=white, drop shadow={opacity=0.15}, 
%                         minimum height=0.8cm, align=left},
%         arrow/.style={-{Stealth[length=2mm]}, line width=1.2pt}
%     ]
    
%     % Configure Java syntax highlighting
%     \lstset{
%         language=Java,
%         basicstyle=\ttfamily\fontsize{7}{8}\selectfont,
%         keywordstyle=\color{blue!70!black}\bfseries,
%         identifierstyle=\color{black},
%         numbers=none,
%         breaklines=true,
%         xleftmargin=3pt,
%         xrightmargin=3pt,
%         backgroundcolor=\color{gray!5},
%         frame=none,
%         escapechar=@,
%         columns=flexible
%     }
    
%     % SFT: Obfuscated Input → Verbatim Reproduction (Left)
%     \node[codebox, draw=red!50, text width=0.45\columnwidth] (sft) {
%         \textbf{\footnotesize SFT: Cognitive Specialization}\\
%         \textit{\tiny Obfuscated input reproduced verbatim}\\[2pt]
%         \begin{lstlisting}
% public class Main_2 {
%   void solve() {
%     Scanner @\colorbox{red!20}{\texttt{IlIlIllIlIlI}}@ = 
%       new Scanner(System.in);
%     int @\colorbox{red!20}{\texttt{llIlIIlIIl}}@ = 
%       @\colorbox{red!20}{\texttt{IlIlIllIlIlI}}@.nextInt();
%     for (int i = @\colorbox{red!20}{\texttt{llIlIIlIIl}}@;
%          i >= 0; --i)
%       System.out.println(i);
%   }
% }
%         \end{lstlisting}
%         \vspace{2pt}
%         \textcolor{red!70!black}{\tiny\bfseries \ding{55} AST$_{\text{synt}}$ = 0.999 (99\% identifiers retained)}
%     };
    
%     % CFT: Successful Deobfuscation (Right)
%     \node[codebox, draw=green!50!black, text width=0.45\columnwidth, right=0.5cm of sft] (cft) {
%         \textbf{\footnotesize CFT: Emergent Bidirectionality}\\
%         \textit{\tiny Semantic recovery with readable identifiers}\\[2pt]
%         \begin{lstlisting}
% public class @\colorbox{green!20}{\texttt{CountingNumbers}}@ {
%   void @\colorbox{green!20}{\texttt{countDown}}@() {
%     Scanner @\colorbox{green!20}{\texttt{scanner}}@ =
%       new Scanner(System.in);
%     int @\colorbox{green!20}{\texttt{number}}@ = 
%       @\colorbox{green!20}{\texttt{scanner}}@.nextInt();
%     for (int i = @\colorbox{green!20}{\texttt{number}}@;
%          i >= 0; --i)
%       System.out.println(i);
%   }
% }
%         \end{lstlisting}
%         \vspace{2pt}
%         \textcolor{green!60!black}{\tiny\bfseries \checkmark AST$_{\text{synt}}$ = 0.17 (83\% identifiers changed)}
%     };
    
%     \end{tikzpicture}
    
%     \caption{Cognitive Specialization vs. Emergent Bidirectionality.}
%     \label{fig:cognitive_specialization}
% \end{figure}


%%%%%%%%%%%%%%%%%%%%%%%%%%%%%%%%%%%%%%%%%%%%%%%%%%


% \begin{figure}[!ht]
%     \centering
%     \scriptsize
%     \begin{tikzpicture}[
%         node distance=0.5cm,
%         codebox/.style={rectangle, rounded corners=3pt, draw=black!40, line width=0.8pt, 
%                         fill=white, drop shadow={opacity=0.15}, 
%                         minimum height=0.8cm, align=left},
%         arrow/.style={-{Stealth[length=2mm]}, line width=1.2pt}
%     ]
    
%     % Configure Java syntax highlighting
%     \lstset{
%         language=Java,
%         basicstyle=\ttfamily\fontsize{7}{8}\selectfont,
%         keywordstyle=\color{blue!70!black}\bfseries,
%         identifierstyle=\color{black},
%         numbers=none,
%         breaklines=true,
%         xleftmargin=3pt,
%         xrightmargin=3pt,
%         backgroundcolor=\color{gray!5},
%         frame=none,
%         escapechar=@,
%         columns=flexible
%     }
    
%     % SFT: Obfuscated Input → Verbatim Reproduction (Left)
%     \node[codebox, draw=red!50, text width=0.45\columnwidth] (sft) {
%         \textbf{\footnotesize SFT: Cognitive Specialization}\\
%         \textit{\tiny Obfuscated input reproduced verbatim}\\[2pt]
%         \begin{lstlisting}
% public class Main_2 {
%   void solve() {
%     Scanner @\colorbox{red!20}{\texttt{IlIlIllIlIlI}}@ = 
%       new Scanner(System.in);
%     int @\colorbox{red!20}{\texttt{llIlIIlIIl}}@ = 
%       @\colorbox{red!20}{\texttt{IlIlIllIlIlI}}@.nextInt();
%     for (int i = @\colorbox{red!20}{\texttt{llIlIIlIIl}}@;
%          i >= 0; --i)
%       System.out.println(i);}
% }
%         \end{lstlisting}
%         \vspace{2pt}
%         \textcolor{red!70!black}{\tiny\bfseries \ding{55} AST$_{\text{synt}}$ = 0.999 (99\% identifiers retained)}
%     };
    
%     % CFT: Successful Deobfuscation (Right)
%     \node[codebox, draw=green!50!black, text width=0.45\columnwidth, right=0.5cm of sft] (cft) {
%         \textbf{\footnotesize CFT: Emergent Bidirectionality}\\
%         \textit{\tiny Semantic recovery with readable identifiers}\\[2pt]
%         \begin{lstlisting}
% public class @\colorbox{green!20}{\texttt{CountingNumbers}}@ {
%   void @\colorbox{green!20}{\texttt{countDown}}@() {
%     Scanner @\colorbox{green!20}{\texttt{scanner}}@ =
%       new Scanner(System.in);
%     int @\colorbox{green!20}{\texttt{number}}@ = 
%       @\colorbox{green!20}{\texttt{scanner}}@.nextInt();
%     for (int i = @\colorbox{green!20}{\texttt{number}}@;
%          i >= 0; --i)
%       System.out.println(i);}
% }
%         \end{lstlisting}
%         \vspace{2pt}
%         \textcolor{green!60!black}{\tiny\bfseries \checkmark AST$_{\text{synt}}$ = 0.17 (83\% identifiers changed)}
%     };
    
%     \end{tikzpicture}
    
%     \caption{Cognitive Specialization vs. Emergent Bidirectionality.}
%     \label{fig:cognitive_specialization}
% \end{figure}


%%%%%%%%%%%%%%%%%%%%%%%%%%%%%%%%%%%%%%%%%%%%%


\begin{figure}[!ht]
  \centering
  \scriptsize
  \begin{tikzpicture}[
      node distance=0.5cm,
      codebox/.style={rectangle, rounded corners=3pt, draw=black!40, line width=0.8pt, 
                      fill=white, drop shadow={opacity=0.15}, 
                      minimum height=0.8cm, align=left},
      arrow/.style={-{Stealth[length=2mm]}, line width=1.2pt}
  ]
  
  % Configure Java syntax highlighting
  \lstset{
      language=Java,
      basicstyle=\ttfamily\fontsize{7}{8}\selectfont,
      keywordstyle=\color{blue!70!black}\bfseries,
      identifierstyle=\color{black},
      numbers=none,
      breaklines=true,
      xleftmargin=3pt,
      xrightmargin=3pt,
      backgroundcolor=\color{gray!5},
      frame=none,
      escapechar=@,
      columns=flexible
  }
  
  % SFT: Obfuscated Input → Verbatim Reproduction (Left)
  \node[codebox, draw=red!50, text width=0.45\columnwidth] (sft) {
      \textbf{\footnotesize SFT: Inversion Curse}\\
      \textit{\tiny Obfuscated input reproduced verbatim}\\[2pt]
      \begin{lstlisting}
void add() {
  int @\colorbox{red!20}{\texttt{l1ll1}}@ = 10;
  int @\colorbox{red!20}{\texttt{l1ll2}}@ = 20;
  int @\colorbox{red!20}{\texttt{l1ll3}}@ = @\colorbox{red!20}{\texttt{l1ll1}}@ + @\colorbox{red!20}{\texttt{l1ll2}}@;
}
      \end{lstlisting}
%      \vspace{2pt}
%      \textcolor{red!70!black}{\tiny\bfseries \ding{55} AST$_{\text{synt}}$ = 0.999 (99\% identifiers retained)}
  };
  
  % CFT: Successful Deobfuscation (Right)
  \node[codebox, draw=green!50!black, text width=0.45\columnwidth, right=0.5cm of sft] (cft) {
      \textbf{\footnotesize CFT: Emergent Bidirectionality}\\
      \textit{\tiny Semantic recovery with readable identifiers}\\[2pt]
      \begin{lstlisting}
void add() 
{
  int @\colorbox{green!20}{\texttt{num1}}@ = 10;
  int @\colorbox{green!20}{\texttt{num2}}@ = 20;
  int @\colorbox{green!20}{\texttt{sum}}@ = @\colorbox{green!20}{\texttt{num1}}@ + @\colorbox{green!20}{\texttt{num2}}@;
}
      \end{lstlisting}
%      \vspace{2pt}
 %     \textcolor{green!60!black}{\tiny\bfseries \checkmark AST$_{\text{synt}}$ = 0.17 (83\% identifiers changed)}
  };
  
  \end{tikzpicture}
  \caption{Inversion     Curse vs. Emergent Bidirectionality.}
  \label{fig:cognitive_specialization}
  % \vspace{-0.6cm}
\end{figure}